\section{实验设置}
\subsection{数据与评价范围}
采用Abilene的144条流和GEANT处理快照的462条流。
Abilene使用公开的X01--X24记录\cite{abilene_data}；
GEANT使用ARI-LLM发布仓库中固定版本的10,772行CSV\cite{ari_release}。
后者被作为确定的处理后数组使用，不将其行号直接换算为原生时间戳或物理单位。

窗口长度为50槽。以零起始、左闭右开的行区间表示，
Abilene的拟合范围为$[0,33335)$，开发范围为$[33384,37484)$；
GEANT相应为$[0,7023)$和$[7072,8322)$。
每套数据在拟合范围内按等间隔起点取512个固定窗口；GEANT拟合窗口存在重叠。
开发窗口互不重叠，分别为82和25个。
选定的512个拟合窗口用于监督训练，完整拟合段用于统计量和相关邻流图的构造，
当前补全窗口的模型输入只包含可见值和掩码。
这是假设具有离线完整历史数据的监督补全设置。

所有结果均来自开发范围；该范围用于检查点选择，不能视作独立最终测试。
历史后续范围也曾被同一项目的早前研究使用，因此本文仅报告开发性结果。
\subsection{三种固定预算缺失条件}
三种条件都严格保留窗口内20\%的数值。
均匀条件中每条流随机保留10/50个位置；
不均匀条件中，随机选取一半流各保留4个位置，其余流各保留16个位置。
共享缺口条件沿用这一流间预算，为一半流设置共同的连续10槽不可见区间，
并在其余位置重新分配观测，保持每条流的观测数量不变。
未受影响流与不均匀条件使用相同观测位置。
因此条件差异来自观测分配和时间形状，不来自总观测量增加。

\subsection{训练、消融与参考}
两WAN、两变体和种子41001/41002构成八条轨迹。
采用AdamW，学习率$10^{-3}$、权重衰减$10^{-4}$、梯度范数上限1、批量8，
不使用学习率调度；全部目标流参与计算，FP32且关闭TF32。
训练掩码和窗口顺序按轮刷新，三条件近似均衡分配；
同种子配对共享这些输入及共有参数初值。
训练掩码、顺序和开发掩码的固定种子分别为61001、51001和71001。

初始轨迹采用早停，随后在已见开发结果后决定导入同一轨迹的模型、优化器和随机状态，
接续到总计120轮。该固定轮数比较是训练时长诊断，
早停阶段与接续阶段只算一次训练，不能充当额外独立重复。
所有变体沿用同一总轮数，以三条件平均开发NMAE最小的检查点作为结果。

唯一训练消融为去Sync，不增加门控、宽度或损失搜索。
参考包含既有四层SPIN、未加入SPIN上下文的既有Direct+Sync，以及逐流线性插值。
线性插值使用可见位置，边缘取端点值，无观测时取拟合期流均值。
旧可训练参考各只保留一个初始化结果，且训练配方与新模型不同：
旧SPIN的上限为300轮，在Abilene/GEANT分别训练137/300轮；
旧Direct+Sync分别训练53/71轮，上限120轮。
这些结果用于完整系统比较，不作为同预算组件归因。

\subsection{指标与完整推理测量}
对每个条件，先合并该开发范围全部缺失位置$\Omega$，计算
\begin{equation}
\begin{aligned}
\mathrm{NMAE}&=\frac{\sum_{\Omega}|\widehat{x}-x|}{\sum_{\Omega}|x|},\\
\mathrm{NRMSE}&=\sqrt{\frac{\sum_{\Omega}(\widehat{x}-x)^2}{\sum_{\Omega}x^2}}.
\end{aligned}
\label{eq:metrics}
\end{equation}
然后对三条件等权取平均。NMAE始终为主要指标；
NRMSE对较大误差更敏感，用于解释取舍。
新模型跨两个种子报告均值和样本标准差，不据此构造显著性结论。

完整推理在同一NVIDIA H20上逐模型运行，采用FP32、关闭TF32，
对两个数据集各自的前8个拟合窗口测试批量1和8。
新两变体与旧SPIN的内部节点分块为32，旧Direct+Sync沿用自身实现。
四种模型均使用seed41001的最佳检查点；每个条件、批量和模型进行5次暖身与15次计时。
计时包含输入传输、全流预测、输出回传、有限值检查、非负截断和观测值回填，
不包含权重加载、拟合统计量及图准备、掩码生成。
表中批量1耗时为三条件各自整批中位数的平均，
批量8耗时再除以8，表示摊销每窗耗时。
显存为批量8三条件中的最大PyTorch已分配值，含当前常驻模型与共享统计量，
不等于整卡占用。此过程只测成本，没有计算新准确率分数。
